% !TEX program = xelatex
% ============================================================
%  学习笔记模板（LaTeX 版）
%  用法：
%   1) 用 XeLaTeX 编译（Overleaf 里在左上角 Menu -> Compiler 选 XeLaTeX）
%   2) 每个 \section 都是一个「问题」，你只负责回答它
%   3) \begin{quote}...\end{quote} 里是填写提示，写完删掉
%   4) 中文文档必须用 XeLaTeX，不要用 pdfLaTeX
% ============================================================

\documentclass[UTF8, 11pt, a4paper]{ctexart}

% ---------- 宏包 ----------
\usepackage{amsmath, amssymb}          % 数学公式
\usepackage{geometry}                  % 页面边距
\geometry{left=2.5cm, right=2.5cm, top=2.5cm, bottom=2.5cm}
\usepackage{booktabs}                  % 好看的三线表
\usepackage{enumitem}                  % 自定义列表
\usepackage{xcolor}                    % 颜色
\usepackage{graphicx}                  % 插图
\usepackage[colorlinks=true, linkcolor=blue, urlcolor=blue]{hyperref}

% ---------- 封面信息 ----------
\title{学习笔记：这里写主题名}
\author{你的名字}
\date{\today}

\begin{document}

\maketitle

% ---------- 元信息表 ----------
\begin{center}
\begin{tabular}{ll}
\toprule
日期 & 2025-01-01 \\
来源 & 课程 / 书 / 视频（写清名字和链接） \\
用时 & 30 分钟 \\
\bottomrule
\end{tabular}
\end{center}

\vspace{1em}
\hrule
\vspace{1em}

% ============================================================
\section{这段内容一句话讲什么？}

\begin{quote}
\itshape
【填写提示】逼自己只用一句话。写不出来，说明还没懂。\\
例：Markdown 是用几个符号代替鼠标点按钮的写字方式。
\end{quote}

（在这里写你的那一句话）


% ============================================================
\section{它讲了哪几件主要的事？}

\begin{quote}
\itshape
【填写提示】只列骨架，不展开细节。3～5 条最合适。
\end{quote}

\begin{itemize}[leftmargin=*, itemsep=0.3em]
  \item 第一件：
  \item 第二件：
  \item 第三件：
\end{itemize}


% ============================================================
\section{每件事具体怎么说？}

\begin{quote}
\itshape
【填写提示】用自己的话复述，别抄原文。抄了就等于没学。\\
需要写公式时：行内公式用 $E = mc^2$ 这样写，独占一行用下面这样——\\
\begin{verbatim}
\[
  \sum_{i=1}^{n} a_i = \frac{n(n+1)}{2}
\]
\end{verbatim}
\end{quote}

\subsection{第一件事：}

\begin{itemize}[leftmargin=*, itemsep=0.3em]
  \item 关键点：
  \item 例子：
  \item 容易搞错的地方：
\end{itemize}

\subsection{第二件事：}

\begin{itemize}[leftmargin=*, itemsep=0.3em]
  \item 关键点：
  \item 例子：
  \item 容易搞错的地方：
\end{itemize}


% ============================================================
\section{我原来的想法错在哪？（最重要的一节）}

\begin{quote}
\itshape
【填写提示】只写「我原以为……其实……」的对比。\\
这一节决定你有没有真的学到东西。
\end{quote}

\begin{center}
\begin{tabular}{p{0.28\textwidth} p{0.28\textwidth} p{0.32\textwidth}}
\toprule
\textbf{我原以为} & \textbf{实际上} & \textbf{差别在哪} \\
\midrule
                &                &              \\
\addlinespace
                &                &              \\
\bottomrule
\end{tabular}
\end{center}


% ============================================================
\section{还没懂 / 存疑的地方}

\begin{quote}
\itshape
【填写提示】写下来，下次优先解决。留疑问不丢人。
\end{quote}

\begin{itemize}[leftmargin=*, label=$\square$, itemsep=0.3em]
  \item 疑问一：
  \item 疑问二：
\end{itemize}


% ============================================================
\section{下一步做什么？}

\begin{quote}
\itshape
【填写提示】必须是可以动手做的，不是「多复习」这种空话。
\end{quote}

\begin{enumerate}[label=(\arabic*), leftmargin=*, itemsep=0.3em]
  \item 马上做：（今天就做，10 分钟内能开始）
  \item 这周做：
  \item 要查的资料：
\end{enumerate}


% ============================================================
\section*{附录：原文摘录}

\begin{quote}
\itshape
【填写提示】值得逐字记的原文放这里，记得标出处和页码 / 时间点。
\end{quote}

\begin{quote}
「　　　」—— 出处，第 00:00 分钟
\end{quote}

\end{document}
